\chapter{Drehstrom und komplexe Wechselstromrechnung}
\label{chap:drehstrom}

\begin{lernziel}
Nach diesem Kapitel kannst du:
\begin{itemize}
  \item Effektivwert, Scheitelwert, Frequenz und Phasenwinkel unterscheiden,
  \item sinusförmige Spannungen und Ströme als Zeigergrößen darstellen,
  \item Impedanzen von Widerstand, Induktivität und Kapazität bilden,
  \item Wirk-, Blind- und Scheinleistung voneinander unterscheiden,
  \item Leiter- und Stranggrößen im symmetrischen Drehstromsystem umrechnen.
\end{itemize}
\end{lernziel}

\section{Warum wir mit Wechselstromrechnung beginnen}

Netze der öffentlichen Energieversorgung arbeiten überwiegend mit
sinusförmigem Drehstrom. Auch wenn reale Netze Oberschwingungen, Unsymmetrien
und transiente Vorgänge enthalten, ist das symmetrische sinusförmige
Grundmodell der Ausgangspunkt für viele Netzberechnungen.

Die zentrale Vereinfachung lautet: Im eingeschwungenen sinusförmigen Zustand
mit einer festen Frequenz können Spannungen und Ströme durch komplexe
Zeigergrößen und Bauteile durch komplexe Impedanzen beschrieben werden.
Differentialgleichungen werden damit für viele stationäre Berechnungen zu
algebraischen Gleichungen.

\section{Sinusförmige Größen}

Eine sinusförmige Spannung kann als
\begin{equation}
  u(t)=\hat U\cos(\omega t+\varphi_u)
  \label{eq:sinusspannung}
\end{equation}
geschrieben werden. Dabei ist:
\begin{description}[style=nextline]
  \item[$u(t)$] Momentanwert der Spannung,
  \item[$\hat U$] Scheitelwert,
  \item[$\omega$] Kreisfrequenz in \si{\radian\per\second},
  \item[$t$] Zeit,
  \item[$\varphi_u$] Phasenwinkel zum gewählten Zeitbezug.
\end{description}

Frequenz und Kreisfrequenz hängen zusammen:
\begin{equation}
  \omega=2\pi f.
\end{equation}
Für eine sinusförmige Größe gilt zwischen Scheitelwert und Effektivwert:
\begin{equation}
  U=\frac{\hat U}{\sqrt{2}},
  \qquad
  I=\frac{\hat I}{\sqrt{2}}.
\end{equation}
Der Effektivwert ist derjenige Gleichspannungs- beziehungsweise Gleichstromwert,
der an einem ohmschen Widerstand dieselbe mittlere Wärmeleistung erzeugen würde.

\begin{merke}
In Netzangaben wie \SI{400}{\volt} oder \SI{20}{\kilo\volt} ist im üblichen
Sinusmodell der Effektivwert gemeint, sofern nichts anderes angegeben ist.
Scheitelwerte dürfen nicht unbemerkt mit Effektivwerten vermischt werden.
\end{merke}

\section{Komplexe Zeigergrößen}

Wir verwenden im Folgenden die Kosinuskonvention. Eine sinusförmige Größe
\[
x(t)=\hat X\cos(\omega t+\varphi_x)
\]
wird durch den Effektivwertzeiger
\begin{equation}
  \phasor{X}=X e^{\jj\varphi_x}
  =X(\cos\varphi_x+\jj\sin\varphi_x)
\end{equation}
repräsentiert. Der Zeitverlauf lässt sich daraus mit
\[
x(t)=\sqrt{2}\,\operatorname{Re}
\left\{\phasor{X}e^{\jj\omega t}\right\}
\]
rekonstruieren.

Die komplexe Schreibweise trennt Betrag und Phasenlage. Das ist entscheidend,
weil an Induktivitäten und Kapazitäten Strom und Spannung nicht phasengleich
sind.

\subsection{Impedanz}

Im sinusförmigen stationären Zustand ist die Impedanz das Verhältnis aus
Spannungs- und Stromzeiger:
\begin{equation}
  Z=\frac{\phasor{U}}{\phasor{I}}=R+\jj X.
  \label{eq:impedanz}
\end{equation}
Hierbei bezeichnet \(R\) den Wirkwiderstand und \(X\) die Reaktanz.
Der Betrag ist
\[
|Z|=\sqrt{R^2+X^2},
\]
und der Impedanzwinkel beträgt
\[
\varphi_Z=\arg(Z)=\arctan\!\left(\frac{X}{R}\right),
\]
wobei für eine robuste Quadrantenwahl rechnerisch die Funktion
\(\operatorname{atan2}(X,R)\) zu verwenden ist.

Für die drei idealisierten Grundelemente gilt:
\begin{align}
  Z_R &= R,\\
  Z_L &= \jj\omega L,\\
  Z_C &= \frac{1}{\jj\omega C}=-\frac{\jj}{\omega C}.
\end{align}
Die induktive Reaktanz \(X_L=\omega L\) wächst mit der Frequenz.
Der Betrag der kapazitiven Reaktanz \( |X_C|=1/(\omega C) \) sinkt mit
wachsender Frequenz.

\section{Wirk-, Blind- und Scheinleistung}

Für Effektivwertzeiger verwenden wir die komplexe Leistung
\begin{equation}
  \underline S=\phasor U\,\phasor I^{*}=P+\jj Q,
  \label{eq:komplexeleistung}
\end{equation}
wobei \(^{*}\) die komplexe Konjugation bezeichnet. Daraus folgen:
\begin{align}
  P &= UI\cos\varphi && \text{Wirkleistung},\\
  Q &= UI\sin\varphi && \text{Blindleistung},\\
  S &= UI=\sqrt{P^2+Q^2} && \text{Scheinleistung}.
\end{align}
Die Vorzeichenkonvention für \(Q\) muss bei jeder konkreten Aufgabe klar sein.
In der üblichen Verbraucherzählpfeilkonvention ist induktive Blindleistung
einer Last positiv und kapazitive Blindleistung negativ.

\section{Symmetrisches Drehstromsystem}

Ein symmetrisches Drehstromsystem besitzt drei sinusförmige Strangspannungen
gleicher Frequenz und gleichen Effektivwerts, die zeitlich um jeweils
\SI{120}{\degree} verschoben sind. Bei positiver Phasenfolge \(L1,L2,L3\)
können die Strangspannungszeiger gegen den Sternpunkt geschrieben werden als
\begin{align}
  \phasor U_{1N}&=U_\mathrm{Str}\angle 0^\circ,\\
  \phasor U_{2N}&=U_\mathrm{Str}\angle -120^\circ,\\
  \phasor U_{3N}&=U_\mathrm{Str}\angle +120^\circ.
\end{align}
Ihre vektorielle Summe ist null:
\[
\phasor U_{1N}+\phasor U_{2N}+\phasor U_{3N}=0.
\]

Für ein symmetrisches System gilt zwischen Leiter-Leiter-Spannung und
Strangspannung:
\begin{equation}
  U_\mathrm{LL}=\sqrt{3}\,U_\mathrm{Str}.
  \label{eq:leiter_strang_spannung}
\end{equation}
Daher entsprechen bei einem üblichen \SI{400}{\volt}-Netz ungefähr
\SI{230}{\volt} zwischen Außenleiter und Neutralleiter. Bei einem
\SI{20}{\kilo\volt}-Netz beträgt die Strangspannung gegen den geerdeten
Sternpunkt im idealisierten Betrieb etwa \SI{11.55}{\kilo\volt}.

\begin{merke}
Die Beziehung \(U_\mathrm{LL}=\sqrt{3}U_\mathrm{Str}\) gilt für das
symmetrische Drehstromsystem. Bei Erdschluss und Sternpunktverschiebung darf
man die tatsächlichen Leiter-Erde-Spannungen nicht unkritisch aus dieser
Normalbetriebsbeziehung ableiten.
\end{merke}

\section{Drehstromleistung}

Für ein symmetrisches Drehstromsystem gilt:
\begin{align}
  P &= \sqrt{3}\,U_\mathrm{LL} I_\mathrm{L}\cos\varphi,\\
  Q &= \sqrt{3}\,U_\mathrm{LL} I_\mathrm{L}\sin\varphi,\\
  S &= \sqrt{3}\,U_\mathrm{LL} I_\mathrm{L}.
\end{align}
Dabei sind \(U_\mathrm{LL}\) die Leiter-Leiter-Spannung und \(I_\mathrm{L}\)
der Leiterstrom. Diese Formeln beziehen sich auf einen symmetrischen Zustand.

\begin{rechenbeispiel}
Ein symmetrischer Verbraucher wird an einem \SI{400}{\volt}-Netz betrieben.
Der Leiterstrom beträgt \SI{32}{\ampere}, der Leistungsfaktor ist
\(\cos\varphi=0{,}90\). Die aufgenommene Wirkleistung ist
\begin{align*}
  P &= \sqrt{3}\cdot \SI{400}{\volt}\cdot
       \SI{32}{\ampere}\cdot 0{,}90\\
    &\approx \SI{19.95}{\kilo\watt}.
\end{align*}
Die Scheinleistung beträgt
\[
S=\sqrt{3}\cdot\SI{400}{\volt}\cdot\SI{32}{\ampere}
 \approx \SI{22.17}{\kilo\volt\ampere}.
\]
Die Blindleistung ergibt sich aus \(Q=\sqrt{S^2-P^2}\), sofern der
induktive beziehungsweise kapazitive Charakter separat bekannt ist:
\[
Q\approx \SI{9.65}{\kilo var}
\]
bei induktivem Verbraucher und der hier angenommenen Vorzeichenkonvention.
\end{rechenbeispiel}

\section{Typische Fehler}

\begin{itemize}
  \item Leiter-Leiter-Spannung und Strangspannung verwechseln.
  \item Effektivwert und Scheitelwert vermischen.
  \item Bei komplexen Größen nur Beträge rechnen und Phasenwinkel ignorieren.
  \item Blindleistung mit Wirkleistung addieren, als wären beide dieselbe Art
        von Größe.
  \item Die Drehstromleistungsformeln auf unsymmetrische Betriebszustände
        anwenden, ohne die Voraussetzungen zu prüfen.
\end{itemize}

\begin{aufgaben}
  \item Eine sinusförmige Spannung hat den Scheitelwert
        \(\hat U=\SI{325}{\volt}\). Bestimme Effektivwert und Frequenz,
        wenn die Periodendauer \SI{20}{\milli\second} beträgt.
  \item Bestimme \(Z\), Betrag und Winkel für \(R=\SI{8}{\ohm}\) und
        \(L=\SI{25}{\milli\henry}\) bei \(f=\SI{50}{\hertz}\).
  \item Ein symmetrischer Verbraucher nimmt bei \SI{400}{\volt} und
        \SI{25}{\ampere} einen Leistungsfaktor von \(0{,}82\) auf.
        Berechne \(P\), \(S\) und den Betrag von \(Q\).
  \item Erkläre, weshalb die Leiter-Erde-Spannung eines gesunden Leiters
        in einem isolierten Mittelspannungsnetz bei einem Erdschluss
        deutlich steigen kann, obwohl die Leiter-Leiter-Spannungen
        näherungsweise unverändert bleiben.
\end{aufgaben}

\section{Kurzlösungen}

\begin{enumerate}[label=\arabic*.]
  \item \(U=\hat U/\sqrt{2}\approx\SI{230}{\volt}\);
        \(f=1/T=\SI{50}{\hertz}\).
  \item \(\omega=2\pi f\approx\SI{314.16}{\radian\per\second}\),
        \(X_L=\omega L\approx\SI{7.85}{\ohm}\),
        \(Z=8+\jj7.85\ \si{\ohm}\),
        \(|Z|\approx\SI{11.21}{\ohm}\), Winkel etwa \(44.5^\circ\).
  \item \(S=\sqrt{3}UI\approx\SI{17.32}{\kilo\volt\ampere}\);
        \(P=S\cdot0.82\approx\SI{14.20}{\kilo\watt}\);
        \(|Q|=S\sqrt{1-0.82^2}\approx\SI{9.92}{\kilo var}\).
  \item Beim Erdschluss verschiebt sich der Sternpunkt gegen Erde. Die
        Leiter-Leiter-Spannungen werden durch die Verkettung der
        Leiterpotentiale bestimmt und können annähernd gleich bleiben,
        während die Leiter-Erde-Spannungen der gesunden Phasen steigen.
\end{enumerate}
